\section{Final identifications}\label{sec:final-identifications}

MWW Theorem~3.10 identifies the iterated quotient with the module after the
three-handles, and Proposition~3.4 gives the four-handle isomorphism
\cite{ManolescuWalkerWedrich2023}.  Propositions~\ref{thm:P0discharge},
\ref{thm:Cdischarge}, and~\ref{thm:Sdischarge} record the status of
Hypotheses~\ref{hyp:P0}--\ref{hyp:P2}; all three are \Open.  For the Johnson
replacement four-handle picture $X_J$, attaching the three $3$-handles
along the belt spheres of the reversed $1$-handles restores the boundary
$S^3$ (Hypothesis~\ref{hyp:P3}), a PL $4$-ball is attached, and the quantum
degree-$494$ summand of the standard sphere module vanishes by the
computation in the status remark of Hypothesis~\ref{hyp:P3} together with
MWW Corollary~3.5.  The staged Kirby pipeline would identify $X_J$ with
$\Sigma_A^0$ once Hypothesis~\ref{hyp:P0} is established.

\begin{remark}[Formalization boundary]\label{rem:lean-boundary}
A Lean development \cite{Lean4} proves Theorem~\ref{thm:A} as a conditional
implication from structures \texttt{ExternalGeometry} and
\texttt{CSExternalGeometry}.  It checks the matrix arithmetic, the value
$2624$, the degree $494$, and the abstract quotient lemmas, but does not
construct the geometric instances proved in this paper.  Appendix~\ref{app:lean-fields}
lists the correspondence with the Lean fields.
\end{remark}
